% related-work.tex -- Cross-Enterprise Verification Paper (RU-V7)

\section{Related Work}

\subsection{ZK Proof Aggregation Systems}

Proof aggregation reduces L1 verification cost by combining multiple proofs
into a single verification operation. SnarkPack~\cite{gailly2022snarkpack}
uses an inner product argument to aggregate Groth16 proofs, achieving 163~ms
verification time for 8,192 proofs with a crossover point at 32 proofs where
aggregation becomes cheaper than individual verification.
aPlonK~\cite{zapico2022aplonk} extends this to PLONK proofs via
multi-polynomial commitment, but requires approximately 300 proofs before
efficiency gains materialize. SnarkFold~\cite{ozdemir2023snarkfold} applies
folding-based aggregation to PLONK, achieving 4.5~ms verification at 4,096
proofs. Nebra UPA~\cite{nebra2024upa} introduces universal proof aggregation
that works across different circuit types, with gas savings starting from just
2 proofs (per-proof cost: $100K/M + 20K$ submission, $350K/N + 7K$
aggregation, plus $25K$ query).

These systems target public rollups and assume all proofs can be collected by
a central aggregator. Our work extends aggregation to the enterprise setting
where each proof originates from a private state tree and the aggregator must
not learn the underlying data.

\subsection{Multi-Chain and Cross-Rollup Architectures}

Polygon AggLayer~\cite{polygon2025agglayer} implements shared bridge
infrastructure with multi-chain ZK proof aggregation, achieving
\$0.002--0.003 per transaction. zkSync Gateway~\cite{zksync2025gateway}
serves as a proof aggregation hub for ZKsync chains at approximately
\$0.0047 per transaction with batch sizes around 3,895 transactions. Both
systems focus on cost reduction through aggregation but operate with public
or semi-public state---enterprise data confidentiality is not a design goal.

StarkPack~\cite{starkpack2024} proposes heterogeneous proof aggregation
(aggregating proofs from different circuits), which is relevant to
cross-enterprise settings where different enterprises may run different
circuit configurations. However, StarkPack does not address the privacy
dimension of cross-enterprise interactions.

\subsection{Enterprise Privacy Systems}

Table~\ref{tab:enterprise-privacy} summarizes production and research systems
addressing enterprise privacy on blockchain.

\begin{table}[t]
\centering
\caption{Enterprise Privacy Systems Comparison}
\label{tab:enterprise-privacy}
\begin{tabular}{@{}lllll@{}}
\toprule
System & Privacy Model & Cross-Ent. & Proof & Status \\
\midrule
Rayls~\cite{rayls2025privacy} & Subnets + ZKP + HE & Hub chain & Enygma & Production \\
Prividium~\cite{zksync2025prividium} & Private execution & Selective & PLONK & Announced \\
zkCross~\cite{zkcross2024} & Cross-chain audit & Compliance & zk-SNARK & Research \\
ZK-InterChain~\cite{zkinterchain2025} & Cross-chain privacy & Limited & Various & Research \\
\textbf{Ours} & Per-ent. validium & Hub-spoke & Groth16 & Research \\
\bottomrule
\end{tabular}
\end{table}

Rayls, developed by Parfin for the Drex (Brazilian CBDC) ecosystem and
adopted by institutions including Nuclea and Cielo, implements private subnets
where each enterprise maintains isolated state with ZK proofs and homomorphic
encryption (the Enygma protocol). Cross-enterprise interactions pass through
a hub chain that coordinates transfers. While Rayls achieves strong
intra-enterprise privacy, the hub chain serves as a coordination point that
sees transaction metadata. ZKsync Prividium~\cite{zksync2025prividium},
announced in December 2025, proposes privacy-default execution with selective
disclosure but has not yet published cross-enterprise verification benchmarks.

\subsection{Cross-Privacy Research}

zkCross~\cite{zkcross2024} introduces privacy-preserving cross-chain auditing
using ZK proofs, targeting compliance scenarios where an auditor verifies
properties across chains without seeing the data. Chainlink
CCIP~\cite{chainlink2024ccip} provides cross-chain interoperability with
privacy features but focuses on message passing rather than state
verification.

\subsection{Benchmark Comparison}

Table~\ref{tab:benchmarks} positions our cross-enterprise verification gas
costs against published benchmarks.

\begin{table}[t]
\centering
\caption{Verification Gas Cost Comparison}
\label{tab:benchmarks}
\begin{tabular}{@{}llr@{}}
\toprule
System & Configuration & Gas/Proof \\
\midrule
Groth16 individual~\cite{nebra2024upa} & 4 public inputs & 205,600 \\
Nebra UPA (N=4)~\cite{nebra2024upa} & Aggregated & 167,457 \\
Nebra UPA (N=32)~\cite{nebra2024upa} & Aggregated & 65,901 \\
\midrule
Ours (Sequential) & 2 ent., 1 int. & 269K/proof \\
Ours (Batched) & 2 ent., 1 int. & 122K/proof \\
Ours (Hub Agg.) & 2 ent., 1 int. & 221K/proof \\
\bottomrule
\end{tabular}
\end{table}

Our per-proof gas costs are higher than Nebra UPA at scale because our system
includes cross-reference verification and storage operations alongside
enterprise proof verification. At equivalent proof counts (N=3), our batched
approach achieves 156K gas per proof including cross-reference storage,
comparable to Nebra UPA at N=4 (167K) which excludes application-level
storage.
